% New mathematical proof component for PA194AX11GH.
% Adapts with attribution the supplied AX11G explicit-binder compiler and
% the current incoming AX11H cache-only compiler and arithmetic templates.
% Requires amsmath, amssymb, amsthm, booktabs.
\section{AX11GH and AX11GHS: exact commands and preservation in iPRA}
\label{sec:ax11gh-preservation}

The two new sources combine the four-operand, eleven-entry AX protocol with
explicit-binder generalization and a cache-only Shift halt test.  This
section proves its source-to-command identity and its refutation compiler
in intuitionistic primitive recursive arithmetic.  It is a written proof
with executable finite binding certificates; it is not a new Lean kernel
build.  The calibration of the numerical calculus to PA is the separately
identified theorem of the current incoming AX11H packet.

\subsection{Numerical calculus and exact command effects}

Use Cantor pairing $p(a,b)=(a+b)(a+b+1)/2+b$ and the constructors
\[
 I(A,B)=4p(A,B),\quad M(a,b)=4p(a,b)+1,\quad
 U(a,B)=4p(B,a)+2,\quad E(a,b)=4p(a,b)+3.
\]
Truth is code zero and the refutation target is code one.
Set $\neg B=I(B,1)$, $B\wedge C=\neg I(B,\neg C)$, and
$\exists_a B=\neg U(a,\neg B)$.  The retained numerical calculus
$\mathcal C_4^Q$ has truth, numerical modus ponens, theorem generalization,
and the following axiom schemes:
\begin{align*}
 h_1(D,H)&=I(\exists_D U(D,H),H),\\
 h_2(D,A,H)&=I(U(D,U(A,H)),U(A,U(D,H))),\\
 h_3(D,A,H)&=I(M(H,D),I(E(H,A),M(A,D))),\\
 h_4(D,A,H)&=I(E(H,D),I(M(A,H),M(A,D))),\\
 h_5(D,H)&=\exists_D E(H,D),\\
 h_6^Q(D,A,H)&=I(U(D,I(A,H)),I(U(D,A),U(D,H))),\\
 \rho(D,A,H,u)&=I(I(I(H,u),D),I(I(D,H),I(A,H))),\quad u\le D,\\
 \Gamma(A)&=I(A,U(A,A)).
\end{align*}
The guarded head is defined by
\begin{align*}
 P(a,b,c)&=I(c,M(c,b)),& V(a,b,c)&=I(M(a,b),U(c,M(a,b))),\\
 S(a,b,c)&=U\bigl(b,I(M(a,b),P(a,b,c))\wedge I(M(c,b),c)\bigr),\\
 Q(a,b,c)&=U(c,I(S(a,b,c),\neg P(a,b,c))),\\
 g(a,b,c)&=\begin{cases}\neg I(V(a,b,c),Q(a,b,c)),&c\ne a,b,\\
                         0,&\text{otherwise.}\end{cases}
\end{align*}
All rules are rules on the exact numerical codes, including the alias
$I(0,0)=0$.  The proof does not presume a disjoint truth constructor.

Write $[a;w]=p(a,w)$ for the numerical stack, and use its total
projections, with $\operatorname{hd}(0)=\operatorname{tl}(0)=0$.
A command store is $(t,w,d,h)$: top, stack, numerical offset, and cache.
For operands $c,x,y,z$, put $D=d+x$ and $H=h+y$.  In top-first order let
\[
\begin{array}{c|l}
q_0&g(H,D,c)\\q_1&\Gamma(c)\\q_2&h_6^Q(D,c,H)\\q_3&h_5(D,H)\\
q_4&h_4(D,c,H)\\q_5&h_3(D,c,H)\\q_6&h_2(D,c,H)\\
q_7&h_1(D,H)\\q_8&\rho(D,c,H,d)\\q_9&U(c,h).
\end{array}
\]
The new commands have exactly the following effects:
\begin{align*}
 \mathsf A_{c,x,y,z}(t,w,d,h)
   &=(0,[q_0,q_1,\ldots,q_9,t;w],z,0),\\
 \mathsf M_x(t,w,d,h)
   &=\begin{cases}(d+x,w,d+x,0),&I(t,d+x)=h,\\
                   (0,w,d+x,0),&I(t,d+x)\ne h,\end{cases}\\
 \mathsf S(t,w,d,h)
   &=\begin{cases}\operatorname{halt},&h=1,\\
       (h,\operatorname{tl}(w),d,\operatorname{hd}(w)),&h\ne1.
       \end{cases}
                                                        \tag{GH}
\end{align*}
AX really has four operands.  Its final push clears $t$ after saving the
guarded head; the guard has already cleared $d,h$.  The following complete
MP reader consumes $z$.  It receives $(t,d,h)=(0,0,0)$, so for $z=0$
its comparison succeeds with result zero and for $z>0$ it fails with
result zero.  In both cases the offset becomes $z$ and the stack is
preserved.  Soundness covers every $z$; the proof compiler selects $z=0$.

\subsection{Complete binding to the supplied primitive source}

The independent script \texttt{check\_ax11gh\_source.py} imports no
Turing-machine simulator or construction compiler.  It expands handwritten
primitive templates for clearing, transfer, duplication, Cantor unpairing,
pairing, multiplication by four, destructive comparison, and fetch.
It checks every primitive instruction and every outgoing edge against the
3,671-site AX11GH source and the 3,668-site AX11GHS source.  Separately it expands the named source tree and
resolves its jumps by subtree intervals.  Both results equal the delivered
instruction array.  The template decompositions contain 947 and 945 instances, respectively.

The ports are fetch 25, decode 34, AX 50, full MP reader 3572,
MP comparison 3642, successful MP copy 3652, return jump 3659, and
Shift 3660.  The source array and all ports agree with the proved AX11H
arithmetic templates.  The changed AX token is the second token: CopyA
replaces CopyD before the first constructor.  Consequently the first saved
universal is $U(c,h)$ instead of $U(d,h)$.  Both copy templates have the
same primitive length and satisfy the same scratch-register conditions.

A free-algebra execution checks all 190 AX tokens, both guard comparisons,
every intermediate scratch precondition, the complete eleven-entry stack,
and all nine final registers.  In particular its residual program register
is $p(z,\ell)$ before the extra MP reader.  Thus the result includes the
fourth field and its exact consumption, rather than silently borrowing the
old three-field interpretation.  The equation checked before that reader is
\[
\begin{split}
 &(n,t,w,c,d,h,0,0,p(y,p(x,p(z,\ell))))\\
 &\quad\longmapsto
 (n,0,[q_0,\ldots,q_9,t;w],0,0,0,0,0,p(z,\ell)).
\end{split}\tag{AX-end}
\]
Here the register order is $(N,T,W,A,D,H,X,Y,L)$.
For each source the checker also replays 882 local exact-duration examples
and 324 composed Shift/MP examples.  These examples are arithmetic regression checks; the
uniform template proof is by the counter inductions described next.

\paragraph{The symmetric companion is independently bound.}
In AX11GHS, all primitive instructions with indices below 3666 coincide
with AX11GH.  At index 3666 the three-instruction restoration loop
\texttt{dec X; inc H; jump} is replaced by \texttt{inc H}, followed by
the jump to root comparison.  More explicitly, after clearing $T$, the
common duplicate loop sends
\[
 (T,H,X)=(0,h,0)\longmapsto(h,0,h).
\]
AX11GH restores $H$ by moving $X$ back to $H$ and increments $X$;
its comparison input is $(T,H,X)=(h,h,1)$.  AX11GHS instead increments
$H$, giving $(h,1,h)$.  The destructive comparison halts exactly when
$h=1$ in both cases and otherwise clears both compared registers,
preserving top $h$.  The common stack-pop continuation then supplies
the next cache.  Thus the two sources implement exactly (GH), including
both halting and nonhalting cases for every command store.

The companion's complete instruction array and outgoing edges are checked
against independently generated templates with duplication replacing
copy-and-restore.  It is not accepted merely because an older GHS packet
used the same source transformation.  At the Shift entry, excluding the
common stack-pop or halt continuation, the respective clocks are
\[
\begin{aligned}
 (2t+1)+(7h+2)+2+C(h,1),\\
 (2t+1)+(4h+1)+2+C(1,h).
\end{aligned}
\]
The first two added steps are the final increment and jump.  The displayed
clock implementation chooses the appropriate variant explicitly, and its
324 composed Shift/MP checks are replayed separately for each source.

\paragraph{The clock is a primitive recursive construction.}
The inherited proof \texttt{source\_command\_completion.tex} specifies
all arithmetic templates and their exact clocks.  For completeness, clear,
transfer, and duplicate of a counter $r$ cost $2r+1$, $3r+1$, and $4r+1$;
copying and restoring $r$ costs $7r+2$ when the scratch register is zero.
The pairing block on $(T,X,Y)=(t,a,0)$ costs
\[
 P(t,a)=4t+2+\sum_{j<a+t}(7j+5)
\]
and leaves $(0,0,p(a,t))$.  Multiplication by four costs $6y+1$.
For unpairing define the primitive recursive succession
\[
 B(a,b)=\begin{cases}(a-1,b+1),&a>0,\\(b+1,0),&a=0,\end{cases}
 \qquad (\operatorname{hd}(n),\operatorname{tl}(n))=B^n(0,0).
\]
Put $r(a,b)=5$ for $a>0$ and $r(0,b)=3b+8$.  The unpairing block costs
\[
 U_*(n)=2+\sum_{j<n}r(B^j(0,0))+3\operatorname{tl}(n).
\]
The comparison of $h,x$ clears both operands and costs
\[
 C(h,x)=\begin{cases}3h+2,&h=x,\\3x+2,&h<x,\\2h+x+4,&h>x.\end{cases}
\]
The source fetch either decrements a positive continuation in two steps or
copies fresh candidate $n$ into $L$, increments $N$, and clears scratch in
$7n+4$ steps.  Costs of composed blocks are added at their successive
arithmetic endpoints.  They are explicit total primitive recursive
functions; there is no minimization over an unknown stopping time.

For each counter loop, fix the initial input and induct on the number of
completed rounds.  The invariant consists of current-register equalities,
a partial-sum clock, and an equality for an iterate of the literal step
function.  Bounded guards are primitive recursive Boolean equations.
This is quantifier-free induction in iPRA.  The finite template composition
then proves (AX-end) and the two other commands uniformly.
For Shift, clear $T$, copy and restore $H$, increment the initially zero
$X$, and invoke the root comparison.  It compares the original $h$ to one;
on failure the stack is unpaired into $H$.  Hence its guard is precisely
$h=1$.  A symmetric implementation may compare one to $h$ instead; the
same comparison theorem gives the identical abstract effect (GH), although
its primitive clock can differ.

Let $\Theta(s)$ be the resulting positive exact source clock of one
scheduler command at store $s$.  Primitive recursion gives cumulative
clock $\tau(k)$ for $k$ commands.  Induction on $k$, expressed as an
iterated-source equality, proves the source/scheduler correspondence.  For
the converse direction, a proposed first literal source halt at time $t$
cannot occur inside $t+1$ positive nonhalting command segments.  This is
bounded finite concatenation, so the source and scheduler halt predicates
are iPRA-provably equivalent from the clean seed at source site 25.

\subsection{Stack words: the cache-only guard matters}

Let $A_0=\mathsf A_{0,0,0,0}$ and define
\[
 \mathsf{Push}=A_0\mathsf S^{10}\mathsf M_0,\qquad
 \mathsf{Pop}=A_0\mathsf S^{12},\qquad
 \mathsf{Discard}=\mathsf M_0.
\]
Every $q_0,\ldots,q_8$ has residue zero modulo four; $q_9=U(c,h)$ has
residue two and is at least two.  Direct iteration of (GH) therefore gives
\begin{align*}
 \mathsf A_{c,x,y,0};\mathsf S^{10}(t,w,d,h)
    &=(q_8,[t;w],0,q_9),\\
 \mathsf A_{c,x,y,0};\mathsf S^{11}(t,w,d,h)
    &=(q_9,w,0,t).
\end{align*}
These are nonhalting calculations for every input.  In Push, the final MP
comparison is between an implication-residue code and a universal-residue
code, so it fails and clears the visible top.  Thus, on every store,
\[
 \mathsf{Push}(t,w,d,h)=(0,[t;w],0,0).
\]
The extra twelfth Shift of Pop tests the saved old top.  The exact Pop law is
\[
 \mathsf{Pop}(t,[a;w],d,h)=
 \begin{cases}\operatorname{halt},&t=1,\\(t,w,0,a),&t\ne1.\end{cases} \tag{Pop}
\]
It is essential not to reuse the unconditional Pop law of the sum-guard
parent.  Finally,
\[
 \mathsf{Discard}(t,w,0,h)=(0,w,0,0),\qquad
 \mathsf M_i(0,w,0,0)=(0,w,i,0).
\]
The latter holds because $i=0$ returns zero on success and $i>0$ fails its
comparison with cache zero.  These identities use total Cantor projections
and hold under $[0;0]=0$; they do not regard stacks as free lists.

\subsection{Axiom selectors and generalization}

For $0\le j\le8$, with the offset adjustment for a nonzero fourth
$\rho$ argument treated below, use
\[
 \mathsf{Sel}_j=
 \mathsf{Discard};\mathsf A_{c,D,H,0};
       \mathsf S^{j+2};\mathsf{Pop}^{9-j}.                 \tag{Sel}
\]
From every $(t,w,0,0)$ it returns $(q_j,w,0,0)$.
After the initial Shifts the selected top is $q_j$, the cache is the next
head, and the remaining generated heads and saved zero lie above $w$.
Every subsequent Pop preserves this selected residue-zero top and removes
one remaining cell.  Hence none of these Pops halts.  The gate selector
is used for a permitted guarded instance or for truth when the guard fails.
For $\rho(D,c,H,u)$ with $u\le D$, the exact selector is
\[
 \mathsf{Discard};\mathsf M_u;
 \mathsf A_{c,D-u,H,0};\mathsf S^{10};\mathsf{Pop}.
\]
The equality $u+(D-u)=D$ ensures the intended dominated instance.

Generalization is particularly short:
\[
 \mathsf{Gen}_i=
 \mathsf{Push};\mathsf S;\mathsf A_{i,0,0,0};\mathsf S^{11}.  \tag{Gen}
\]
From $(a,w,0,0)$, Push and the next Shift give $(0,w,0,a)$; that Shift
sees zero in the cache before moving $a$ there.  AX saves $U(i,a)$ and
clears the visible cache, so the eleven following Shifts return
$(U(i,a),w,0,0)$ without halting.  This calculation remains valid for
$a=1$.  It requires no offset-setting command because the first AX
operand is the binder.

\subsection{Proof compilation on arbitrary protected stacks}

At an MP node, suppose the major proof concludes $I(a,b)$ and the minor
proof concludes $a$.  Define the command word recursively by
\[
\begin{split}
 W_{\mathrm{MP}(e,f)}={}&\mathsf{Discard};\mathsf{Push};
 W_e;\mathsf{Push};W_f;\mathsf{Push};\mathsf S;\mathsf S;
                         \mathsf M_b;\mathsf{Pop}.
\end{split}\tag{MP-word}
\]
The first Discard and Push reserve a zero under the two premises.
If the recursive words return, the store before the two Shifts is
$(0,[a,I(a,b),0;w],0,0)$.  The first Shift is safe; the second halts
if $a=1$.  Otherwise the result is $(a,[0;w],0,I(a,b))$.
MP gives $(b,[0;w],b,0)$.  By (Pop), the final Pop halts if $b=1$
and otherwise returns $(b,w,0,0)$.  The zero reservation is necessary
because successful MP no longer pushes a dummy truth cell.

Use Discard for a truth leaf, (Sel) for an axiom leaf, and append (Gen)
at a generalization node.  The precise inductive claim is:
\[
 W_e(t,w,0,0)\text{ either halts or returns }
       (\operatorname{concl}(e),w,0,0)                  \tag{Compiler}
\]
for every valid finite numerical derivation $e$ and every $t,w$.
Both possibilities are allowed even when a subproof concludes one.
No consistency hypothesis is used to exclude an early halt.

\paragraph{A bounded induction formula.}
One must not perform iPRA induction directly on an unbounded
``for all stores'' assertion.  Fix a checked proof table, a height bound
$F$ (leaves have height one), and the requested input $t,w$.
Let $K\ge1$ bound the conclusion code of every node, and set
\[
 R_0=\max(1,K,t,w),\qquad R_{j+1}=(K+R_j+3)^4.
\]
The only recursively requested MP stacks, starting with stack $y$, are
$[0;y]$ and $[I(a,b),0;y]$.  For $y\le R$ and $I(a,b)\le K$, the
bound $p(v,x)\le(v+x+1)^2$ places both below $(K+R+3)^4$.
The recursively requested tops are zero.  Generalization requests its
premise at the unchanged store.

At stage $k\le F$, take the finite conjunction of (Compiler) for every
proof-table node of height at most $k$ and every $x,y\le R_{F-k}$.
Execution of a finite command word is primitive recursive; thus this
finite conjunction is one primitive recursive Boolean equation.
At stage $k+1$, recursive calls made from inputs bounded by
$R_{F-k-1}$ lie within the preceding allowance $R_{F-k}$.
The local word identities prove the remaining cases, including all
one/nonone tests by decidable arithmetic.  Quantifier-free induction
proves the final desired instance.  The originally fixed proof and store
parameters are generalized only afterwards.  All conjunctions and their
bounds are explicit primitive recursive functions.

There is also a primitive recursive word-length bound: a selector has
at most 122 commands, a generalization suffix has 25, and an MP node
adds 53 commands to its two subwords.  Hence
$L_0=0$, $L_{j+1}=2L_j+128$, and $L_j=128(2^j-1)$ are safe bounds.
These length bounds do not claim a small physical running time.

\subsection{Every carried store admits a halting refutation word}

A plain Shift need not detect top one: from $(1,w,0,0)$ it does not halt.
The finishing word
\[
 \mathsf{Finish}=\mathsf{Push};\mathsf S;\mathsf S
\]
halts from every store with top one.  Push gives
$(0,[1;w],0,0)$, the first Shift gives $(0,w,0,1)$, and the second halts.
Alternatively Pop detects a top one.  For a proof $e$ of one, therefore,
\[
                 A_0;W_e;\mathsf{Finish}                \tag{Ref-word}
\]
halts from every incoming store.  The initial AX normalizes the offset
and cache; the compiler either halts or returns top one with its protected
stack.  The carried stack is never assumed to reset between candidates.

The full command-word encoding is
\begin{align*}
 e(\epsilon)&=0,& e(a::v)&=1+b(a,e(v)),\\
 b(\mathsf S,r)&=p(0,r),&
 b(\mathsf M_x,r)&=p(1,p(x,r)),\\
 b(\mathsf A_{c,x,y,z},r)&=p(c+2,p(y,p(x,p(z,r)))).
\end{align*}
A fresh candidate omits the first outer one.  The decoded AX fields occur
in order $y,x,z$.  The symbolic checker verifies all inverse identities,
including zero aliases, and explicitly checks that the old three-field
encoding consumes the wrong continuation.

For any suffix (Ref-word), prefixing $\mathsf A_{k,0,0,0}$ gives fresh
candidate code
\[
       p(k+2,p(0,p(0,p(0,e(\text{Ref-word}))))).
\]
These candidates are cofinal.  Each prefix terminates and hands an
arbitrary store to the universally halting suffix.  Every positive
continuation strictly decreases because it is decoded from a code one
smaller than itself, and every extracted Cantor component is at most its
pair code.  Thus every nonhalting fresh candidate finishes in a bounded
number of command steps.  A bounded loop through the finitely many earlier
candidate numbers computes a primitive recursive halting witness for the
chosen cofinal candidate unless an earlier halt occurs.  Its invariant is
again a primitive recursive Boolean equation.  This establishes scheduler
completeness in iPRA without an unbounded stopping-time argument.

\subsection{Soundness and preservation of the consistency calibration}

Annotate top, cache, and every live canonical Cantor-stack entry with a
$\mathcal C_4^Q$ proof of its numerical code.  The initial store has only
zeros, with their truth proof.  AX inserts the displayed axioms, the old
top, and the generalization $U(c,h)$ of its annotated cache.
The dominated condition for $\rho(D,c,H,d)$ is $d\le d+x$.
Both guard branches are covered; a failed gate yields truth.  Its final
extra MP reader puts truth in both visible positions for every $z$.
Successful MP combines the two premise proofs, failure yields truth,
and Shift moves existing annotations.  If Shift halts, its annotated
cache is exactly code one.  The zero-stack collapse is covered by the
truth proof at virtual zero entries.  Proof construction, concatenation,
and stack annotation are primitive recursive, and their finite-prefix
correctness is a bounded Boolean induction.

Consequently each new literal source scheduler $R$ satisfies
\[
 \mathrm{iPRA}\vdash
 H_{R}
 \ \longleftrightarrow\
 \exists e\,\operatorname{Prf}_{\mathcal C_4^Q}(e,1).     \tag{Numerical}
\]
Positive exact source-to-tape simulations and the exact blank startup
transfer (Numerical) to each displayed binary table.  This transfer uses
the frame and arithmetic-box proof dependencies documented in the main
report; a bounded run sample alone would not supply it.
The current incoming QAX11H construction has the same numerical-calculus
refutation predicate, with its separately proved offset-binder compiler.
Thus each new table and that exact reference table have iPRA-provably
equivalent blank-start halting predicates.

The current incoming AX11H packet supplies a written calibration of this
predicate to the repository's native PA-refutation sentence, using its
model-uniform $\mathrm I\Sigma_1$ proof compilers and the recorded transfer
to iPRA.  Applying that theorem gives
\[
 \boxed{\mathrm{iPRA}\vdash
       \operatorname{Con}(\mathrm{PA})
       \longleftrightarrow\neg H_M,
       \qquad M\in\{\mathrm{AX11GH},\mathrm{AX11GHS}\}.}
\]
The new relative preservation and command compiler do not by themselves
reprove that extensive PA calibration.  The latter is an inherited written
mathematical theorem, not a consistency assumption and not a completed
end-to-end kernel certificate.  The conclusion is obtained by negating the
existential refutation equivalence intuitionistically; no Markov principle
or double-negation elimination for arbitrary arithmetical assertions is
introduced.
